% GENERATED FILE. Edit canonical lessons, then run npm run generate:books.
% canonical-source-hash: fnv1a64:65966ea4653a084b
% canonical-lessons: MR-C231-ekatra, MR-C231-kadacit, MR-C231-sivay, MR-C231-arthat, MR-C231-mhanje, MR-R231-first-pass-words-from-chapters-223-226, MR-R231-second-pass-words-from-chapters-227-231

\chapter{Together, Perhaps, Besides, Of course, That is to say}
\label{ch:mr-together-perhaps-besides-of-course-that-is-to-say}
\hlchaptermodality{231}

\begin{chapteropening}
\textbf{By the end of this chapter:} \emph{I can say together, perhaps, besides, of course and that is to say.}

Complete the last lesson of chapter 231: I can say together, perhaps, besides, of course and that is to say.
\end{chapteropening}

\section[\texorpdfstring{\mr{एकत्र} (ekatra)}{एकत्र (ekatra)}]{\mr{एकत्र} (ekatra) — together}

\label{lesson:MR-C231-ekatra}



\begin{quote}
Before the new one: say the Marathi for meanwhile, then the Marathi for for good.
\end{quote}

\subsection*{You'll want to know: \mr{एकत्र}}
\textbf{\mr{एकत्र}} — \emph{ekatra} — \textquotedblleft{}together\textquotedblright{}.

\begin{grammarlens}[title={What you've built}]
One more word for this chapter.
\end{grammarlens}

\subsection*{Your turn}
\begin{itemize}
\raggedright
  \item \emph{Say it:} \emph{ekatra}
  \item \emph{Say it:} \emph{ekatra}, once more
  \item \emph{Say it:} say \emph{ekatra}
  \item \emph{Recall:} say the Marathi for meanwhile, then the Marathi for for good, then say \emph{ekatra} again
\end{itemize}

\begin{tcolorbox}[breakable,colback=teal!4,colframe=teal!35!black,title={Before you move on}]
What does \mr{एकत्र} mean? (\textquotedblleft{}Together\textquotedblright{}.) Say it once more.
\end{tcolorbox}
\section[\texorpdfstring{\mr{कदाचित} (kadācit)}{कदाचित (kadācit)}]{\mr{कदाचित} (kadācit) — perhaps}

\label{lesson:MR-C231-kadacit}



\begin{quote}
Before the new one: say the Marathi for for good, then the Marathi for together.
\end{quote}

\subsection*{You'll want to know: \mr{कदाचित}}
\textbf{\mr{कदाचित}} — \emph{kadācit} — \textquotedblleft{}perhaps\textquotedblright{}.

\begin{grammarlens}[title={What you've built}]
One more word for this chapter.
\end{grammarlens}

\subsection*{Your turn}
\begin{itemize}
\raggedright
  \item \emph{Say it:} \emph{kadācit}
  \item \emph{Say it:} \emph{kadācit}, once more
  \item \emph{Say it:} say \emph{kadācit}
  \item \emph{Recall:} say the Marathi for for good, then the Marathi for together, then say \emph{kadācit} again
\end{itemize}

\begin{tcolorbox}[breakable,colback=teal!4,colframe=teal!35!black,title={Before you move on}]
What does \mr{कदाचित} mean? (\textquotedblleft{}Perhaps\textquotedblright{}.) Say it once more.
\end{tcolorbox}
\section[\texorpdfstring{\mr{शिवाय} (śivāy)}{शिवाय (śivāy)}]{\mr{शिवाय} (śivāy) — besides, apart from}

\label{lesson:MR-C231-sivay}



\begin{quote}
Before the new one: say the Marathi for together, then the Marathi for perhaps.
\end{quote}

\subsection*{You'll want to know: \mr{शिवाय}}
\textbf{\mr{शिवाय}} — \emph{śivāy} — \textquotedblleft{}besides, apart from\textquotedblright{}.

\begin{grammarlens}[title={What you've built}]
One more word for this chapter.
\end{grammarlens}

\subsection*{Your turn}
\begin{itemize}
\raggedright
  \item \emph{Say it:} \emph{śivāy}
  \item \emph{Say it:} \emph{śivāy}, once more
  \item \emph{Say it:} say \emph{śivāy}
  \item \emph{Recall:} say the Marathi for together, then the Marathi for perhaps, then say \emph{śivāy} again
\end{itemize}

\begin{tcolorbox}[breakable,colback=teal!4,colframe=teal!35!black,title={Before you move on}]
What does \mr{शिवाय} mean? (\textquotedblleft{}Besides, apart from\textquotedblright{}.) Say it once more.
\end{tcolorbox}
\section[\texorpdfstring{\mr{अर्थात} (arthāt)}{अर्थात (arthāt)}]{\mr{अर्थात} (arthāt) — of course}

\label{lesson:MR-C231-arthat}



\begin{quote}
Before the new one: say the Marathi for perhaps, then the Marathi for besides.
\end{quote}

\subsection*{You'll want to know: \mr{अर्थात}}
\textbf{\mr{अर्थात}} — \emph{arthāt} — \textquotedblleft{}of course\textquotedblright{}.

\begin{grammarlens}[title={What you've built}]
One more word for this chapter.
\end{grammarlens}

\subsection*{Your turn}
\begin{itemize}
\raggedright
  \item \emph{Say it:} \emph{arthāt}
  \item \emph{Say it:} \emph{arthāt}, once more
  \item \emph{Say it:} say \emph{arthāt}
  \item \emph{Recall:} say the Marathi for perhaps, then the Marathi for besides, then say \emph{arthāt} again
\end{itemize}

\begin{tcolorbox}[breakable,colback=teal!4,colframe=teal!35!black,title={Before you move on}]
What does \mr{अर्थात} mean? (\textquotedblleft{}Of course\textquotedblright{}.) Say it once more.
\end{tcolorbox}
\section[\texorpdfstring{\mr{म्हणजे} (mhaṇje)}{म्हणजे (mhaṇje)}]{\mr{म्हणजे} (mhaṇje) — that is to say}

\label{lesson:MR-C231-mhanje}



\begin{quote}
Before the new one: say the Marathi for besides, then the Marathi for of course.
\end{quote}

\subsection*{You'll want to know: \mr{म्हणजे}}
\textbf{\mr{म्हणजे}} — \emph{mhaṇje} — \textquotedblleft{}that is to say\textquotedblright{}.

\begin{grammarlens}[title={What you've built}]
One more word for this chapter.
\end{grammarlens}

\subsection*{Your turn}
\begin{itemize}
\raggedright
  \item \emph{Say it:} \emph{mhaṇje}
  \item \emph{Say it:} \emph{mhaṇje}, once more
  \item \emph{Say it:} say \emph{mhaṇje}
  \item \emph{Recall:} say the Marathi for besides, then the Marathi for of course, then say \emph{mhaṇje} again
\end{itemize}

\begin{tcolorbox}[breakable,colback=teal!4,colframe=teal!35!black,title={Before you move on}]
What does \mr{म्हणजे} mean? (\textquotedblleft{}That is to say\textquotedblright{}.) Say it once more.
\end{tcolorbox}
\section[(dialogue)]{Words from chapters 223-226 — a first pass}

\label{lesson:MR-R231-first-pass-words-from-chapters-223-226}



\begin{quote}
Say the words for of course and that is to say.
\end{quote}

\subsection*{Your turn}
Say these aloud:

\begin{itemize}
\raggedright
  \item to repair, to take a photo, to explain, to translate, to make a decision
  \item to grow old, to get divorced, to move house, to rest, to get bored
  \item to convince, to criticise, to prove, to depend, to avoid
  \item to suggest, to give back, to refuse, to take a risk, to solve
\end{itemize}

\begin{tcolorbox}[breakable,colback=teal!4,colframe=teal!35!black,title={Before you move on}]
To repair? (\textbf{\mr{दुरुस्त} \mr{करणे}}, \emph{durust karṇe}.) To take a photo? (\textbf{\mr{फोटो} \mr{काढणे}}, \emph{phoṭo kāḍhṇe}.) To explain? (\textbf{\mr{समजावणे}}, \emph{samjāvṇe}.) To translate? (\textbf{\mr{भाषांतर} \mr{करणे}}, \emph{bhāṣāntar karṇe}.) To make a decision? (\textbf{\mr{निर्णय} \mr{घेणे}}, \emph{nirṇay gheṇe}.) To grow old? (\textbf{\mr{म्हातारे} \mr{होणे}}, \emph{mhātāre hoṇe}.) To get divorced? (\textbf{\mr{घटस्फोट} \mr{घेणे}}, \emph{ghaṭasphoṭ gheṇe}.) To move house? (\textbf{\mr{घर} \mr{बदलणे}}, \emph{ghar badalṇe}.) To rest? (\textbf{\mr{आराम} \mr{करणे}}, \emph{ārām karṇe}.) To get bored? (\textbf{\mr{कंटाळणे}}, \emph{kaṇṭāḷṇe}.) To convince? (\textbf{\mr{पटवणे}}, \emph{paṭavṇe}.) To criticise? (\textbf{\mr{टीका} \mr{करणे}}, \emph{ṭīkā karṇe}.) To prove? (\textbf{\mr{सिद्ध} \mr{करणे}}, \emph{siddha karṇe}.) To depend? (\textbf{\mr{अवलंबून} \mr{असणे}}, \emph{avlambūn asṇe}.) To avoid? (\textbf{\mr{टाळणे}}, \emph{ṭāḷṇe}.) To suggest? (\textbf{\mr{सुचवणे}}, \emph{sucavṇe}.) To give back? (\textbf{\mr{परत} \mr{देणे}}, \emph{parat deṇe}.) To refuse? (\textbf{\mr{नाकारणे}}, \emph{nākārṇe}.) To take a risk? (\textbf{\mr{धोका} \mr{पत्करणे}}, \emph{dhokā patkarṇe}.) To solve? (\textbf{\mr{सोडवणे}}, \emph{soḍavṇe}.)
\end{tcolorbox}
\section[(dialogue)]{Words from chapters 227-231 — a second pass}

\label{lesson:MR-R231-second-pass-words-from-chapters-227-231}



\begin{quote}
Say the words for to destroy and to exaggerate.
\end{quote}

\subsection*{Your turn}
Say these aloud:

\begin{itemize}
\raggedright
  \item to destroy, to exaggerate, to celebrate, to float, to chase
  \item to praise, to go forward, to go back, to defend, to disappear
  \item a tablet, an ointment, a bandage; a strip, an ambulance, a diet
  \item the beginning, a deadline, daily, meanwhile, for good
  \item together, perhaps, besides, of course, that is to say
\end{itemize}

\begin{tcolorbox}[breakable,colback=teal!4,colframe=teal!35!black,title={Before you move on}]
To destroy? (\textbf{\mr{नष्ट} \mr{करणे}}, \emph{naṣṭa karṇe}.) To exaggerate? (\textbf{\mr{अतिशयोक्ती} \mr{करणे}}, \emph{atiśayoktī karṇe}.) To celebrate? (\textbf{\mr{साजरा} \mr{करणे}}, \emph{sājrā karṇe}.) To float? (\textbf{\mr{तरंगणे}}, \emph{taraṅgṇe}.) To chase? (\textbf{\mr{पाठलाग} \mr{करणे}}, \emph{pāṭhlāg karṇe}.) To praise? (\textbf{\mr{कौतुक} \mr{करणे}}, \emph{kautuk karṇe}.) To go forward? (\textbf{\mr{पुढे} \mr{जाणे}}, \emph{puḍhe jāṇe}.) To go back? (\textbf{\mr{मागे} \mr{जाणे}}, \emph{māge jāṇe}.) To defend? (\textbf{\mr{बचाव} \mr{करणे}}, \emph{bacāv karṇe}.) To disappear? (\textbf{\mr{नाहीसे} \mr{होणे}}, \emph{nāhīse hoṇe}.) A tablet? (\textbf{\mr{गोळी}}, \emph{goḷī}.) An ointment? (\textbf{\mr{मलम}}, \emph{malam}.) A bandage; a strip? (\textbf{\mr{पट्टी}}, \emph{paṭṭī}.) An ambulance? (\textbf{\mr{रुग्णवाहिका}}, \emph{rugṇavāhikā}.) A diet? (\textbf{\mr{आहार}}, \emph{āhār}.) The beginning? (\textbf{\mr{सुरुवात}}, \emph{suruvāt}.) A deadline? (\textbf{\mr{मुदत}}, \emph{mudat}.) Daily? (\textbf{\mr{दररोज}}, \emph{darroj}.) Meanwhile? (\textbf{\mr{दरम्यान}}, \emph{darmyān}.) For good? (\textbf{\mr{कायम}}, \emph{kāyam}.) Together? (\textbf{\mr{एकत्र}}, \emph{ekatra}.) Perhaps? (\textbf{\mr{कदाचित}}, \emph{kadācit}.) Besides? (\textbf{\mr{शिवाय}}, \emph{śivāy}.) Of course? (\textbf{\mr{अर्थात}}, \emph{arthāt}.) That is to say? (\textbf{\mr{म्हणजे}}, \emph{mhaṇje}.)
\end{tcolorbox}
